\chapter{最適化への応用}
    \section{\texorpdfstring{$\norm{A\bm{x} + \bm{b}}_2^2$}{||Ax+b||\_2}の最小化条件}
        \subsection{制約なしの場合}
            \begin{shadebox}
                $A\in\complexNumbers^{m\times n},\;\bm{b} \in \complexNumbers^m, \bm{x} \in \complexNumbers^n$とする。
                $f(\bm{x}) \coloneqq \norm{A\bm{x} + \bm{b}}_2^2$がある$\mathring{\bm{x}} \in \complexNumbers^n$で最小となる必要十分条件は次式である。
                \[ \HerConj{A}A\mathring{\bm{x}} = -\HerConj{A}\bm{b} \]
                特に$A$が列フルランクであれば$\HerConj{A}A$は正則(\ref{列フルランク行列を掛けあわせて正則行列を作る})なので$\mathring{\bm{x}}$は次式で一意に定まる。
                \[ \mathring{\bm{x}} = -(\HerConj{A}A)^{-1}\HerConj{A}\bm{b} \]
            \end{shadebox}
            \begin{proof}
                \quad\par
                \ref{複素引数凸関数の例}より$f$は凸であるので、$\mathring{\bm{x}}$が$f$を最小化する必要十分条件は次式である。
                \[ f(\mathring{\bm{x}} + \mathrm{d}\bm{x}) - f(\mathring{\bm{x}}) = o(\|\mathrm{d}\bm{x}\|) \; (\mathrm{d}\bm{x} \to \bm{0}) \]
                これは$f(\mathring{\bm{x}} + \mathrm{d}\bm{x})$の$\mathrm{d}\bm{x}$に関する1次の項が0であることと同値である。
                この1次の項を$g(\mathrm{d}\bm{x})$とすると
                \begin{align*}
                    g(\mathrm{d}\bm{x}) &= \HerConj{(A\mathrm{d}\bm{x})}(A\mathring{\bm{x}} + \bm{b}) + \HerConj{(A\mathring{\bm{x}} + \bm{b})}(A\mathrm{d}\bm{x}) = 2\Re{(\HerConj{\mathring{\bm{x}}}\HerConj{A}A + \HerConj{\bm{b}}A)\mathrm{d}\bm{x}}
                \end{align*}
                であるから、$\mathring{\bm{x}}$に関する条件は$\HerConj{\mathring{\bm{x}}}\HerConj{A}A + \HerConj{\bm{b}}A = O$、すなわち$\HerConj{A}A\mathring{\bm{x}} = -\HerConj{A}\bm{b}$である。
            \end{proof}
        \subsection{線形制約付きの場合}
            \begin{shadebox}
                $f$の定義は前小節のものとする。
                $m\in\naturalNumbers,\;\bm{c}_i \in \complexNumbers^n, d_i\in\complexNumbers,\;g_i(\bm{x}) \coloneqq \bm{c}_i^\top\bm{x}+d_i$とする。
                ある$\mathring{\bm{x}} \in \complexNumbers^n, \mathring{\bm{\lambda}} \in\complexNumbers^m$に対して$g_i(\mathring{\bm{x}}) = 0\;(i=1,\dotsc,m),\; \HerConj{(A\mathring{\bm{x}}+\bm{b})}A = \sum_{i=1}^m \mathring{\lambda}_i\bm{c}_i^\top$であるとき、$f$は制約条件$g_i(\bm{x})=0\;(i=1,\dotsc,m)$の下で$\mathring{\bm{x}}$に於いて最小値をとる。
            \end{shadebox}
            \begin{proof}
                \quad\par
                $\bm{h} \in \complexNumbers^n,\;g_i(\mathring{\bm{x}} + \bm{h}) = 0\;(i=1,\dotsc,m)$とする。
                $g_i(\mathring{\bm{x}} + \bm{h}) = 0 \iff \bm{c}_i^\top\bm{h} = 0$に注意する。
                \begin{align*}
                    f(\mathring{\bm{x}} + \bm{h}) &= f(\mathring{\bm{x}}) + 2\Re{\HerConj{(A\mathring{\bm{x}}+\bm{b})}A\bm{h}} + \HerConj{\bm{h}}\HerConj{A}A\bm{h} = f(\mathring{\bm{x}}) + 2\Re{\sum_{i=1}^m \mathring{\lambda}_i\bm{c}_i^\top\bm{h}} + \HerConj{\bm{h}}\HerConj{A}A\bm{h} \\
                    &= f(\mathring{\bm{x}}) + \HerConj{\bm{h}}\HerConj{A}A\bm{h} \geq f(\mathring{\bm{x}})
                \end{align*}
            \end{proof}